\documentclass[main.tex]{subfiles}
% Eigene Register: niemals die kanonischen B08-Artefakte überschreiben.
% Die gefilterten Importdateien erzeugt scripts/csb-pilot.py aus Band 08.
\begin{luacode*}
  local directory = "registry/csb/"
  local b08_import = "b08-external"
  if tex.jobname == "_B08-csb-reading" then b08_import = "b08-reading" end
  thmlookup.registry_path = directory .. tex.jobname .. ".registry.tsv"
  thmlookup.debug_path = directory .. tex.jobname .. ".debug.log"
  thmlookup.prepare_run()
  for _, predecessor in ipairs(banddeps.graph.B08.predecessors) do
    local base = banddeps.graph[predecessor].artifact_base
    tex.sprint("\\externaldocument{" .. base .. "}[../_" .. predecessor .. ".pdf]")
  end
  tex.sprint("\\externaldocument{" .. directory .. b08_import .. "}[../_B08.pdf]")
  tex.sprint("\\externaldocument{registry/csb/b11-application}[../_B11.pdf]")
  banddeps.load_predecessor_registries("B08")
  thmlookup.load_registry_file(directory .. b08_import .. ".registry.tsv")
  thmlookup.load_registry_file("registry/csb/b11-application.registry.tsv")
\end{luacode*}
\renewcommand{\theHtable}{\FormulaBandID.\arabic{chapter}.\arabic{table}}
\renewcommand{\theHfigure}{\FormulaBandID.\arabic{chapter}.\arabic{figure}}

\CSBReferenzfalse
\CSBLesetextfalse
\CSBTabellentrue
\title{Cantor--Schröder--Bernstein\\[.5em]\large Beweistabellen zu Band 08}
\author{Martin Kunze}
\date{}
\hypersetup{pdftitle={Band 08 - Cantor-Bernstein - Beweistabellen},pdfauthor={Martin Kunze}}
\begin{document}
\hypersetup{pageanchor=false}
\maketitle
\hypersetup{pageanchor=true,hypertexnames=false}
\hypertarget{csb.proofs}{}
\chapter*{Zur Benutzung}
Dieser Auszug enthält die Definition, die acht Sätze und sämtliche
tabellarischen Ableitungen des Cantor--Schröder--Bernstein-Abschnitts
aus Band~08. Die drei mit H bezeichneten Hilfsresultate bleiben mit
ihren eigenen Verweiszielen erhalten. Die Nummern entsprechen dem
Fachband. Frühere Ergebnisse werden dort oder in den Vorgängerbänden
nachgeschlagen.

Die zugehörige \CSBReadingLink{} erklärt die Konstruktion und ihre
entscheidenden Beweisschritte in zusammenhängender mathematischer Sprache.
Die Tabellen sind explizite Manuskriptableitungen; die technischen
Verweisprüfungen ersetzen keine maschinelle Prüfung der Schlussregeln.
\input{registry/csb/position.tex}
\ProofWideReasonsNearFormula
\section{Der Satz von Cantor--Schröder--Bernstein}
% Gemeinsame Deklarationen: unveränderte Aussagen, Kontexte und IDs.
\newcommand{\CSBPartDef}{%
\FormulaDefDeltaK[Cantor--Bernstein-Teilmenge]{%
  \begin{aligned}[t]
    &F\colon A\inj B\dsep G\colon B\inj A\vdash{}\\[-2pt]
    &\CBPart{F}{G}\coloneqq
      \{x\in A\mid \forall X\in\powerset(A)\,\bigl(\\[-2pt]
    &\qquad ((A\setminus G[B])\subseteq X\land G[F[X]]\subseteq X)
      \rightarrow x\in X\bigr)\}
  \end{aligned}%
}{CantorBernsteinPartDef}{%
  \DeltaRow{Mengen}{A\dsep B\dsep X\dsep x}
  \DeltaRow{Funktionen}{F\dsep G}
  \DeltaRow{\textbf{Neues Symbol}}{}
  \DeltaRow{Cantor--Bernstein-Teilmenge}{\CBPart{F}{G}}
}
}

\newcommand{\CSBPartSubset}{%
\FormulaThmDeltaK[Die Cantor--Bernstein-Teilmenge liegt im Ausgangsbereich]{%
  F\colon A\inj B\dsep G\colon B\inj A
  \vdash \CBPart{F}{G}\subseteq A%
}{CantorBernsteinPartSubset}{%
  \DeltaRow{Mengen}{A\dsep B\dsep x}
  \DeltaRow{Funktionen}{F\dsep G}
}
}

\newcommand{\CSBPartContainsSeed}{%
\FormulaThmDeltaK[Der Startbereich liegt im Cantor--Bernstein-Teil]{%
  F\colon A\inj B\dsep G\colon B\inj A
  \vdash A\setminus G[B]\subseteq\CBPart{F}{G}%
}{CantorBernsteinPartContainsSeed}{%
  \DeltaRow{Mengen}{A\dsep B\dsep X\dsep x}
  \DeltaRow{Funktionen}{F\dsep G}
}
}

\newcommand{\CSBPartMinimal}{%
\FormulaThmDeltaK[Minimalität des Cantor--Bernstein-Teils]{%
  \begin{aligned}[t]
    &F\colon A\inj B\dsep G\colon B\inj A\dsep X\subseteq A\dsep
      A\setminus G[B]\subseteq X\dsep G[F[X]]\subseteq X\vdash{}\\[-2pt]
    &\CBPart{F}{G}\subseteq X
  \end{aligned}%
}{CantorBernsteinPartMinimal}{%
  \DeltaRow{Mengen}{A\dsep B\dsep X\dsep Y\dsep x\dsep z}
  \DeltaRow{Funktionen}{F\dsep G}
}
}

\newcommand{\CSBPartClosed}{%
\FormulaThmDeltaK[Abgeschlossenheit des Cantor--Bernstein-Teils]{%
  F\colon A\inj B\dsep G\colon B\inj A
  \vdash G[F[\CBPart{F}{G}]]\subseteq\CBPart{F}{G}%
}{CantorBernsteinPartClosed}{%
  \DeltaRow{Mengen}{A\dsep B\dsep X\dsep y}
  \DeltaRow{Funktionen}{F\dsep G}
}
}

\newcommand{\CSBPartFixedPoint}{%
\FormulaThmDeltaK[Fixpunktgleichung des Cantor--Bernstein-Teils]{%
  \begin{aligned}[t]
    F\colon A\inj B\dsep G\colon B\inj A\vdash
    \CBPart{F}{G}={}&(A\setminus G[B])\\[-2pt]
      &{}\cup G[F[\CBPart{F}{G}]]
  \end{aligned}%
}{CantorBernsteinPartFixedPoint}{%
  \DeltaRow{Mengen}{A\dsep B}
  \DeltaRow{Funktionen}{F\dsep G}
}
}

\newcommand{\CSBComplementIdentity}{%
\FormulaThmDeltaK[Komplementform der Fixpunktgleichung]{%
  F\colon A\inj B\dsep G\colon B\inj A
  \vdash G[B]\setminus G[F[\CBPart{F}{G}]]=A\setminus\CBPart{F}{G}%
}{CantorBernsteinComplementIdentity}{%
  \DeltaRow{Mengen}{A\dsep B\dsep y}
  \DeltaRow{Funktionen}{F\dsep G}
}
}

\newcommand{\CSBSecondBranchImage}{%
\FormulaThmDeltaK[Bild des zweiten Cantor--Bernstein-Zweigs]{%
  \begin{aligned}[t]
    F\colon A\inj B\dsep G\colon B\inj A\vdash
    G[B\setminus F[\CBPart{F}{G}]]
      =A\setminus\CBPart{F}{G}
  \end{aligned}%
}{CantorBernsteinSecondBranchImage}{%
  \DeltaRow{Mengen}{A\dsep B}
  \DeltaRow{Funktionen}{F\dsep G}
}
}

\newcommand{\CSBFixedInjections}{%
\FormulaThmDeltaK[Cantor--Schröder--Bernstein für feste Injektionen]{%
  F\colon A\inj B\dsep G\colon B\inj A
  \vdash \exists H\colon A\bij B%
}{CantorBernsteinFixedInjections}{%
  \DeltaRow{Mengen}{A\dsep B}
  \DeltaRow{Funktionen}{F\dsep G\dsep H}
}
}


\ifCSBReferenz
  Sind Injektionen \(F\colon A\inj B\) und \(G\colon B\inj A\)
gegeben, so gibt es eine Bijektion zwischen \(A\) und \(B\).
Die folgenden Aussagen beschreiben die dazu verwendete Teilmenge und
die beiden bijektiven Zweige.

Die vollständigen Beweise dieses Abschnitts stehen in zwei verknüpften
Begleitfassungen: Die \CSBReadingLink{} entwickelt die Konstruktion
in einem zusammenhängenden mathematischen Beweis; die \CSBProofsLink{}
führen die einzelnen Ableitungsschritte aus. Beide verwenden die
hier angegebenen Satznummern.

\Needspace{14\baselineskip}
\CSBPartDef
\Needspace{14\baselineskip}
\CSBPartSubset
\Needspace{14\baselineskip}
\CSBPartContainsSeed
\Needspace{14\baselineskip}
\CSBPartMinimal
\Needspace{16\baselineskip}
\CSBPartClosed
\begingroup
  \ProofPartSetParent{\theformulaThm}
  \setcounter{proofpartnr}{0}
  \CSBAuxReference{CantorBernsteinPartClosedUniversal}
\endgroup
\Needspace{14\baselineskip}
\CSBPartFixedPoint
\Needspace{16\baselineskip}
\CSBComplementIdentity
\begingroup
  \ProofPartSetParent{\theformulaThm}
  \setcounter{proofpartnr}{0}
  \CSBAuxReference{CantorBernsteinComplementForward}
  \CSBAuxReference{CantorBernsteinComplementBackward}
\endgroup
\Needspace{14\baselineskip}
\CSBSecondBranchImage
\Needspace{14\baselineskip}
\CSBFixedInjections
\noindent\emph{Beweis:} \CSBReadingLink{} und \CSBProofsLink{}.

\else
\ifCSBLesetext
  Der Satz von Cantor--Schröder--Bernstein besagt: Gibt es Injektionen
von \(A\) nach \(B\) und von \(B\) nach \(A\), so gibt es eine
Bijektion zwischen beiden Mengen. Umgekehrt liefern eine Bijektion und
ihre Umkehrfunktion sofort Injektionen in beide Richtungen.

\CSBHeading{Die Idee der Konstruktion}
Seien \(F\colon A\inj B\) und \(G\colon B\inj A\) die gegebenen
Injektionen. Wir suchen eine Teilmenge \(C\subseteq A\), für die folgende
beiden Einschränkungen Bijektionen sind:
\[
  \begin{array}{c@{\qquad}c@{\qquad}c}
    \text{Teilmenge von }A && \text{Teilmenge von }B\\[4pt]
    C & \xrightarrow{\ F\ } & F[C]\\[5pt]
    A\setminus C & \xleftarrow{\ G\ } & B\setminus F[C].
  \end{array}
\]
Die obere Zuordnung benutzt \(F\). Die untere Zuordnung kehren wir um:
Jedem \(a\in A\setminus C\) soll das eindeutige
\(b\in B\setminus F[C]\) mit \(G(b)=a\) zugeordnet werden.
So werden die beiden disjunkten Teile von \(A\) bijektiv auf die beiden
disjunkten Teile von \(B\) abgebildet. Zusammen ergeben sie die gesuchte
Bijektion zwischen \(A\) und \(B\).

Die obere Zuordnung ist für jede Teilmenge \(C\subseteq A\) bijektiv
auf \(F[C]\), weil \(F\) injektiv ist. Bei der unteren Zuordnung
ist die Injektivität ebenfalls gegeben. Entscheidend ist also die Gleichung
\[
  G[B\setminus F[C]]=A\setminus C.
\]
Sie stellt sicher, dass jeder Punkt des verbleibenden Teils von \(A\)
genau ein Urbild im verbleibenden Teil von \(B\) besitzt.

Der folgende Beweis konstruiert dafür die kleinste Teilmenge
\(C\subseteq A\), die \(A\setminus G[B]\) enthält und für die
\(G[F[C]]\subseteq C\) gilt. Aus ihrer Minimalität folgt
\[
  C=(A\setminus G[B])\cup G[F[C]].
\]
Diese Fixpunktgleichung liefert die benötigte Bildgleichung und damit
die beiden Bijektionen des Schemas.

  % Zusammenhängender Prosabeweis; Aussagen und Kennungen bleiben gemeinsam.
\Needspace{6\baselineskip}
\begin{proof}[Beweis]
Seien \(F\colon A\inj B\) und \(G\colon B\inj A\) gegeben.
Wir konstruieren eine Bijektion, indem wir zwei Einschränkungen dieser
Funktionen passend zusammensetzen. Die Buchstaben \(S,C,D\) und
\(\mathcal K\) bezeichnen dabei nur Hilfsmengen dieses Beweises.

\medskip\noindent
\textbf{Die kleinste abgeschlossene Teilmenge.}
Setze
\[
  \begin{aligned}
    S&:=A\setminus G[B],\\
    \mathcal K&:=\{X\in\powerset(A)\mid
      S\subseteq X\land G[F[X]]\subseteq X\}.
  \end{aligned}
\]
Die Familie \(\mathcal K\) ist durch Aussonderung aus der Potenzmenge
von \(A\) eine Menge. Sie ist nicht leer: Da \(F[A]\subseteq B\)
und \(G[B]\subseteq A\) gelten, erfüllt \(A\) beide Bedingungen und
gehört selbst zu \(\mathcal K\).

Wir definieren \(C\) als den Durchschnitt dieser Familie:
\[
  C=\{x\in A\mid \forall X\in\mathcal K\;x\in X\}
   =\bigcap\mathcal K.
\]
Der Durchschnitt ist hier über eine nichtleere Mengenfamilie gebildet.
Da \(A\in\mathcal K\) gilt, ist er eine Teilmenge von \(A\). Jedes Mitglied von
\(\mathcal K\) enthält \(S\); folglich gilt \(S\subseteq C\).
Sei nun \(X\in\mathcal K\) beliebig. Aus \(C\subseteq X\)
folgt durch Monotonie der Bildbildung
\(G[F[C]]\subseteq G[F[X]]\). Weil \(X\) zu \(\mathcal K\)
gehört, gilt nach deren Definition \(G[F[X]]\subseteq X\).
Zusammen ergibt dies
\[
  G[F[C]]\subseteq G[F[X]]\subseteq X.
\]
Jeder Punkt aus \(G[F[C]]\) liegt somit in allen Mitgliedern von
\(\mathcal K\), also in \(C\). Damit ist \(C\) selbst ein Mitglied
von \(\mathcal K\). Zugleich ist \(C\) eine Teilmenge jedes anderen Mitglieds
und ist daher die kleinste Teilmenge von \(A\), welche die beiden
geforderten Bedingungen erfüllt.

\medskip\noindent
\textbf{Die Fixpunktgleichung.}
Aus der Minimalität ergibt sich die stärkere Aussage
\[
  C=S\cup G[F[C]].
\]
Setze dazu \(D:=S\cup G[F[C]]\). Die gerade bewiesenen Inklusionen
liefern \(D\subseteq C\), also insbesondere \(D\subseteq A\).
Wiederum wegen der Monotonie der Bildbildung gilt
\[
  G[F[D]]\subseteq G[F[C]]\subseteq D.
\]
Auch \(S\subseteq D\) gilt nach Definition. Somit gehört \(D\)
zu \(\mathcal K\), und die Minimalität von \(C\) liefert
\(C\subseteq D\). Zusammen mit \(D\subseteq C\) ist die
Fixpunktgleichung bewiesen. Sie beschreibt \(C\) als Vereinigung
von \(S\) und dem Bild von \(C\) unter \(G\circ F\).

\medskip\noindent
\textbf{Das Bild des zweiten Zweiges.}
Die Fixpunktgleichung besagt
\(C=(A\setminus G[B])\cup G[F[C]]\). Beim Übergang zum
Komplement in \(A\) entsteht nach der De-Morganschen Regel zunächst
ein Schnitt:
\[
  \begin{aligned}
    A\setminus C
      &=\bigl(A\setminus(A\setminus G[B])\bigr)
        \cap\bigl(A\setminus G[F[C]]\bigr)\\
      &=G[B]\cap\bigl(A\setminus G[F[C]]\bigr)\\
      &=G[B]\setminus G[F[C]].
  \end{aligned}
\]
Die zweite Zeile benutzt \(G[B]\subseteq A\): Ein Punkt von \(A\),
der nicht in \(A\setminus G[B]\) liegt, liegt gerade in \(G[B]\).
In der letzten Zeile kann die Bedingung, zu \(A\) zu gehören,
entfallen, weil jeder Punkt von \(G[B]\) sie bereits erfüllt.
Es bleiben also genau die Punkte von \(G[B]\), die nicht zu
\(G[F[C]]\) gehören.

Auch \(G[B]\setminus(A\cap G[F[C]])\) ist dieselbe Menge:
Aus \(F[C]\subseteq B\) folgt
\(G[F[C]]\subseteq G[B]\subseteq A\), also
\(A\cap G[F[C]]=G[F[C]]\). Der zusätzliche Schnitt mit \(A\)
ändert die Menge somit nicht.

Mit der Injektivität von \(G\) lässt sich die rechte Seite auch
als Bild einer Differenz schreiben:
\[
  G[B]\setminus G[F[C]]=G[B\setminus F[C]].
\]
Dies lässt sich unmittelbar an den Elementen prüfen. Ist
\(b\in B\setminus F[C]\), so gehört \(G(b)\) zu \(G[B]\).
Läge \(G(b)\) auch in \(G[F[C]]\), gäbe es ein \(c\in F[C]\)
mit \(G(b)=G(c)\). Die Injektivität von \(G\) ergäbe \(b=c\),
im Widerspruch zu \(b\notin F[C]\). Umgekehrt besitzt jeder Punkt
aus \(G[B]\setminus G[F[C]]\) ein Urbild \(b\in B\).
Dieses Urbild kann nicht zu \(F[C]\) gehören und liegt deshalb in
\(B\setminus F[C]\). Insgesamt haben wir damit gezeigt:
\[
  G[B\setminus F[C]]=A\setminus C.
\]

\medskip\noindent
\textbf{Die beiden Bijektionen.}
Eine injektive Funktion ist auf jeder Teilmenge ihres Definitionsbereichs
bijektiv auf ihr Bild. Wegen \(C\subseteq A\) ist daher
\[
  F\restriction_C\colon C\bij F[C]
\]
eine Bijektion. Entsprechend ist die Einschränkung von \(G\) auf
\(B\setminus F[C]\) bijektiv auf dessen Bild. Die gerade bewiesene
Bildgleichung identifiziert dieses Bild mit \(A\setminus C\):
\[
  G\restriction_{B\setminus F[C]}\colon
  B\setminus F[C]\bij A\setminus C.
\]
Erst diese Bijektion kehren wir um und erhalten
\[
  J:=\bigl(G\restriction_{B\setminus F[C]}\bigr)^{-1}
  \colon A\setminus C\bij B\setminus F[C].
\]

\medskip\noindent
\textbf{Zusammensetzen.}
Definiere die Funktion \(H\) auf \(A\) durch
\[
  H(a):=
  \begin{cases}
    F(a),&a\in C,\\
    J(a),&a\in A\setminus C.
  \end{cases}
\]
Die beiden Definitionsbereiche sind disjunkt und vereinigen sich zu
\(A\); jeder Zweig nimmt seine Werte in \(B\) an. Damit ist \(H\)
eine wohldefinierte Funktion von \(A\) nach \(B\), also gerade die
Fallfunktion \(\CaseFun{C}{A\setminus C}{F\restriction_C}{J}\).

Innerhalb jedes Zweiges ist \(H\) injektiv. Auch zwischen den Zweigen
können keine gleichen Werte auftreten: Der erste Zweig hat das Bild
\(F[C]\), der zweite das dazu disjunkte Bild \(B\setminus F[C]\).
Folglich ist \(H\) auf ganz \(A\) injektiv. Jeder Punkt aus \(B\)
gehört zu einem dieser beiden Bildbereiche und wird vom zugehörigen
bijektiven Zweig getroffen. Somit ist \(H\) auch surjektiv und daher
die gesuchte Bijektion \(H\colon A\bij B\).
\end{proof}

\CSBHeading{Einordnung}
Der vorstehende Beweis benötigt weder eine Folge natürlicher Zahlen noch eine
Rekursionskonstruktion oder das Auswahlaxiom. Auch leere Mengen sind
eingeschlossen: Die Familie \(\mathcal K\) bleibt nichtleer, weil
sie \(A\) enthält. Die Hilfsmenge \(C\) dient der Konstruktion;
für spätere Anwendungen genügt das Ergebnis, dass zwei Injektionen
in entgegengesetzten Richtungen die Existenz einer Bijektion sichern.

\fi

\ifCSBTabellen
  \CSBHeading{Aussagen und tabellarische Ableitungen}
  % Beide Ausgabearten führen dieselben Hauptdeklarationen in gleicher Reihenfolge aus.
\CSBPartDef
\CSBProof{%
\Needspace{14\baselineskip}
}

\CSBPartSubset
\CSBProof{%
\begin{tabproof}
  \proofstep{1}{F\colon A\inj B}{\rA}
  \proofstep{2}{G\colon B\inj A}{\rA}
  \proofstep{3}{x\in\CBPart{F}{G}}{\rA}
  \proofstep{1,2,3}{x\in A}{%
    \FormulaRefAuto{CantorBernsteinPartDef}[def]{1,2,3}}
  \proofstep{1,2}{\CBPart{F}{G}\subseteq A}{%
    \FormulaRefAuto{SubsetIntroductionRule}[thm]{[3]\vdots 4}}
\end{tabproof}
}
\Needspace{14\baselineskip}

\CSBPartContainsSeed
\CSBProof{%
\begin{tabproofwide}
  \proofstepwidestar[1]{F\colon A\inj B}{\rA}
  \proofstepwidestar[2]{G\colon B\inj A}{\rA}
  \proofstepwidestar[3]{x\in A\setminus G[B]}{\rA}
  \proofstepwidestar[4]{%
    (A\setminus G[B])\subseteq X\land G[F[X]]\subseteq X}{\rA}
  \proofstepwidestar[4]{A\setminus G[B]\subseteq X}{\rAEa{4}}
  \proofstepwidestar[3,4]{x\in X}{%
    \FormulaRefAuto{A\subseteq B\dsep x\in A\vdash x\in B}{5,3}}
  \proofstepwidestar[3]{%
    ((A\setminus G[B])\subseteq X\land G[F[X]]\subseteq X)
      \rightarrow x\in X}{\rRI{4,6}}
  \proofstepwidestar[3]{%
    X\in\powerset(A)\rightarrow\bigl(
      ((A\setminus G[B])\subseteq X\land G[F[X]]\subseteq X)
        \rightarrow x\in X\bigr)}{%
    \FormulaRefAuto{Q \vdash P \rightarrow Q}[thm]{7}}
  \proofstepwidestar[3]{%
    \forall X\in\powerset(A)\,\bigl(
      ((A\setminus G[B])\subseteq X\land G[F[X]]\subseteq X)
      \rightarrow x\in X\bigr)}{\rUI{8}}
  \proofstepwidestar[3]{x\in A}{%
    \FormulaRefAuto{x\in A\setminus B\vdash x\in A}{3}}
  \proofstepwidestar[1,2,3]{x\in\CBPart{F}{G}}{%
    \FormulaRefAuto{CantorBernsteinPartDef}[def]{1,2,\rAI{10,9}}}
  \proofstepwidestar[1,2]{A\setminus G[B]\subseteq\CBPart{F}{G}}{%
    \FormulaRefAuto{SubsetIntroductionRule}[thm]{[3]\vdots 11}}
\end{tabproofwide}
}
\Needspace{14\baselineskip}

\CSBPartMinimal
\CSBProof{%
\begin{tabproofwide}
  \proofstepwidestar[1]{F\colon A\inj B}{\rA}
  \proofstepwidestar[2]{G\colon B\inj A}{\rA}
  \proofstepwidestar[3]{X\subseteq A}{\rA}
  \proofstepwidestar[4]{A\setminus G[B]\subseteq X}{\rA}
  \proofstepwidestar[5]{G[F[X]]\subseteq X}{\rA}
  \proofstepwidestar[6]{x\in\CBPart{F}{G}}{\rA}
  \proofstepwidestar[3]{X\in\powerset(A)}{%
    \FormulaRefAuto{X\in\powerset(A)\eqvdash X\subseteq A}{3}}
  \proofstepwidestar[4,5]{%
    (A\setminus G[B])\subseteq X\land G[F[X]]\subseteq X}{\rAI{4,5}}
  \proofstepwidestar[1,2]{%
    \CBPart{F}{G}=\bigl\{z\in A\bigm|
      \forall Y\in\powerset(A)\,\bigl(
        ((A\setminus G[B])\subseteq Y\land G[F[Y]]\subseteq Y)
          \rightarrow z\in Y\bigr)\bigr\}}{%
    \FormulaRefAuto{CantorBernsteinPartDef}[def]{1,2}}
  \proofstepwidestar[1,2,6]{%
    \forall Y\in\powerset(A)\,\bigl(
      ((A\setminus G[B])\subseteq Y\land G[F[Y]]\subseteq Y)
        \rightarrow x\in Y\bigr)}{%
    \FormulaRefAuto{x \in A\dsep A=\{x \in B \mid P(x)\}
      \vdash P(x)}[thm]{6,9}}
  \proofstepwidestar[1,2,6]{%
    X\in\powerset(A)\rightarrow\bigl(
      ((A\setminus G[B])\subseteq X\land G[F[X]]\subseteq X)
        \rightarrow x\in X\bigr)}{\rUE{10}}
  \proofstepwidestar[1,2,3,6]{%
    ((A\setminus G[B])\subseteq X\land G[F[X]]\subseteq X)
      \rightarrow x\in X}{\rRE{11,7}}
  \proofstepwidestar[1,2,3,4,5,6]{x\in X}{\rRE{12,8}}
  \proofstepwidestar[1,2,3,4,5]{\CBPart{F}{G}\subseteq X}{%
    \FormulaRefAuto{SubsetIntroductionRule}[thm]{[6]\vdots 13}}
\end{tabproofwide}
}
\Needspace{24\baselineskip}

\CSBPartClosed
\CSBProof{%
\begin{tabproofsplitwide}
  \Needspace{8\baselineskip}
  \hypertarget{csb.helper.CantorBernsteinPartClosedUniversal}{}
  \proofpartwideindR[Bildpunkte liegen in jeder zulässigen Teilmenge]{%
    \begin{aligned}[t]
      &F\colon A\inj B\dsep G\colon B\inj A\dsep
        y\in G[F[\CBPart{F}{G}]]\vdash{}\\[-2pt]
      &\forall X\in\powerset(A)\,\bigl(
        ((A\setminus G[B])\subseteq X\land G[F[X]]\subseteq X)
        \rightarrow y\in X\bigr)
    \end{aligned}%
  }{%
    F\colon A\inj B\dsep G\colon B\inj A\dsep
    y\in G[F[\CBPart{F}{G}]]\vdash
    \forall X\in\powerset(A)\,\bigl(
      ((A\setminus G[B])\subseteq X\land G[F[X]]\subseteq X)
      \rightarrow y\in X\bigr)
  }[CantorBernsteinPartClosedUniversal]
    \proofstepwidestar[1]{F\colon A\inj B}{\rA}
    \proofstepwidestar[2]{G\colon B\inj A}{\rA}
    \proofstepwidestar[3]{y\in G[F[\CBPart{F}{G}]]}{\rA}
    \proofstepwidestar[4]{X\in\powerset(A)}{\rA}
    \proofstepwidestar[5]{%
      (A\setminus G[B])\subseteq X\land G[F[X]]\subseteq X}{\rA}
    \proofstepwidestar[4]{X\subseteq A}{%
      \FormulaRefAuto{X\in\powerset(A)\eqvdash X\subseteq A}{4}}
    \proofstepwidestar[5]{A\setminus G[B]\subseteq X}{\rAEa{5}}
    \proofstepwidestar[5]{G[F[X]]\subseteq X}{\rAEb{5}}
    \proofstepwidestar[1,2,4,5]{\CBPart{F}{G}\subseteq X}{%
      \FormulaRefAuto{CantorBernsteinPartMinimal}[thm]{1,2,6,7,8}}
    \proofstepwidestar[1]{F\colon A\to B}{%
      \FormulaRefAuto{Injektive Funktion}[def]{1}}
    \proofstepwidestar[2]{G\colon B\to A}{%
      \FormulaRefAuto{Injektive Funktion}[def]{2}}
    \proofstepwidestar[1,2,4,5]{F[\CBPart{F}{G}]\subseteq F[X]}{%
      \FormulaRefAuto{F\colon M\to N\dsep A\subseteq B\dsep B\subseteq M
        \vdash F[A]\subseteq F[B]}[thm]{10,9,6}}
    \proofstepwidestar[1,4]{F[X]\subseteq B}{%
      \FormulaRefAuto{C\subseteq A\dsep F\colon A\to B
        \vdash F[C]\subseteq B}[thm]{6,10}}
    \proofstepwidestar[1,2,4,5]{%
      G[F[\CBPart{F}{G}]]\subseteq G[F[X]]}{%
      \FormulaRefAuto{F\colon M\to N\dsep A\subseteq B\dsep B\subseteq M
        \vdash F[A]\subseteq F[B]}[thm]{11,12,13}}
    \proofstepwidestar[1,2,3,4,5]{y\in G[F[X]]}{%
      \FormulaRefAuto{A\subseteq B\dsep x\in A\vdash x\in B}[thm]{14,3}}
    \proofstepwidestar[1,2,3,4,5]{y\in X}{%
      \FormulaRefAuto{A\subseteq B\dsep x\in A\vdash x\in B}[thm]{8,15}}
    \proofstepwidestar[1,2,3,4]{%
      ((A\setminus G[B])\subseteq X\land G[F[X]]\subseteq X)
        \rightarrow y\in X}{\rRI{5,16}}
    \proofstepwidestar[1,2,3]{%
      \forall X\in\powerset(A)\,\bigl(
        ((A\setminus G[B])\subseteq X\land G[F[X]]\subseteq X)
        \rightarrow y\in X\bigr)}{\rUI{\rRI{4,17}}}
  \closeproofpartwideindR

  \Needspace{8\baselineskip}

  \proofpartwide{Zugehörigkeit zu \(A\) und Schluss}
    \proofstepwidestar[1]{F\colon A\inj B}{\rA}
    \proofstepwidestar[2]{G\colon B\inj A}{\rA}
    \proofstepwidestar[3]{y\in G[F[\CBPart{F}{G}]]}{\rA}
    \proofstepwidestar[1,2,3]{%
      \forall X\in\powerset(A)\,\bigl(
        ((A\setminus G[B])\subseteq X\land G[F[X]]\subseteq X)
        \rightarrow y\in X\bigr)}{%
      \FormulaRefAuto{CantorBernsteinPartClosedUniversal}[thm]{1,2,3}}
    \proofstepwidestar[1]{F\colon A\to B}{%
      \FormulaRefAuto{Injektive Funktion}[def]{1}}
    \proofstepwidestar[2]{G\colon B\to A}{%
      \FormulaRefAuto{Injektive Funktion}[def]{2}}
    \proofstepwidestar[1,2]{\CBPart{F}{G}\subseteq A}{%
      \FormulaRefAuto{CantorBernsteinPartSubset}[thm]{1,2}}
    \proofstepwidestar[1,2]{F[\CBPart{F}{G}]\subseteq B}{%
      \FormulaRefAuto{C\subseteq A\dsep F\colon A\to B
        \vdash F[C]\subseteq B}[thm]{7,5}}
    \proofstepwidestar[1,2]{G[F[\CBPart{F}{G}]]\subseteq A}{%
      \FormulaRefAuto{C\subseteq A\dsep F\colon A\to B
        \vdash F[C]\subseteq B}[thm]{8,6}}
    \proofstepwidestar[1,2,3]{y\in A}{%
      \FormulaRefAuto{A\subseteq B\dsep x\in A\vdash x\in B}[thm]{9,3}}
    \proofstepwidestar[1,2,3]{y\in\CBPart{F}{G}}{%
      \FormulaRefAuto{CantorBernsteinPartDef}[def]{1,2,\rAI{10,4}}}
    \proofstepwidestar[1,2]{%
      G[F[\CBPart{F}{G}]]\subseteq\CBPart{F}{G}}{%
      \FormulaRefAuto{SubsetIntroductionRule}[thm]{[3]\vdots 11}}
  \closeproofpartwide
\end{tabproofsplitwide}
}
\Needspace{14\baselineskip}

\CSBPartFixedPoint
\CSBProof{%
\begin{tabproofsplitwide}
\Needspace{8\baselineskip}
\proofpartwide{Erste Inklusion}
\proofstepwidestar[1]{F\colon A\inj B}{\rA}
  \proofstepwidestar[2]{G\colon B\inj A}{\rA}
  \proofstepwidestar[1,2]{A\setminus G[B]\subseteq\CBPart{F}{G}}{%
    \FormulaRefAuto{CantorBernsteinPartContainsSeed}[thm]{1,2}}
  \proofstepwidestar[1,2]{%
    G[F[\CBPart{F}{G}]]\subseteq\CBPart{F}{G}}{%
    \FormulaRefAuto{CantorBernsteinPartClosed}[thm]{1,2}}
  \proofstepwidestar[1,2]{%
    (A\setminus G[B])\cup G[F[\CBPart{F}{G}]]
      \subseteq\CBPart{F}{G}}{%
    \FormulaRefAuto{UnionSubsetCommonSuperset}[thm]{3,4}}

\closeproofpartwide

\Needspace{8\baselineskip}

\proofpartwide{Zweite Inklusion}
  \setcounter{proofstepnr}{5}% Fortlaufende Schritte dieses Beweises.
\proofstepwidestar[]{%
    A\setminus G[B]\subseteq
      (A\setminus G[B])\cup G[F[\CBPart{F}{G}]]}{%
    \FormulaRefAuto{A\subseteq A\cup B}[thm]}
  \proofstepwidestar[1,2]{\CBPart{F}{G}\subseteq A}{%
    \FormulaRefAuto{CantorBernsteinPartSubset}[thm]{1,2}}
  \proofstepwidestar[1,2]{%
    (A\setminus G[B])\cup G[F[\CBPart{F}{G}]]\subseteq A}{%
    \FormulaRefAuto{A\subseteq B\dsep B\subseteq C\vdash A\subseteq C}[thm]{5,7}}
  \proofstepwidestar[1]{F\colon A\to B}{%
    \FormulaRefAuto{Injektive Funktion}[def]{1}}
  \proofstepwidestar[2]{G\colon B\to A}{%
    \FormulaRefAuto{Injektive Funktion}[def]{2}}
  \proofstepwidestar[1,2]{%
    F[(A\setminus G[B])\cup G[F[\CBPart{F}{G}]]]
      \subseteq F[\CBPart{F}{G}]}{%
    \FormulaRefAuto{F\colon M\to N\dsep A\subseteq B\dsep B\subseteq M
      \vdash F[A]\subseteq F[B]}[thm]{9,5,7}}
  \proofstepwidestar[1,2]{F[\CBPart{F}{G}]\subseteq B}{%
    \FormulaRefAuto{C\subseteq A\dsep F\colon A\to B
      \vdash F[C]\subseteq B}[thm]{7,9}}
  \proofstepwidestar[1,2]{%
    G[F[(A\setminus G[B])\cup G[F[\CBPart{F}{G}]]]]
      \subseteq G[F[\CBPart{F}{G}]]}{%
    \FormulaRefAuto{F\colon M\to N\dsep A\subseteq B\dsep B\subseteq M
      \vdash F[A]\subseteq F[B]}[thm]{10,11,12}}
  \proofstepwidestar[]{%
    G[F[\CBPart{F}{G}]]\subseteq
      (A\setminus G[B])\cup G[F[\CBPart{F}{G}]]}{%
    \FormulaRefAuto{B\subseteq A\cup B}[thm]}
  \proofstepwidestar[1,2]{%
    G[F[(A\setminus G[B])\cup G[F[\CBPart{F}{G}]]]]
      \subseteq (A\setminus G[B])\cup G[F[\CBPart{F}{G}]]}{%
    \FormulaRefAuto{A\subseteq B\dsep B\subseteq C\vdash A\subseteq C}[thm]{13,14}}
  \proofstepwidestar[1,2]{%
    \CBPart{F}{G}\subseteq
      (A\setminus G[B])\cup G[F[\CBPart{F}{G}]]}{%
    \FormulaRefAuto{CantorBernsteinPartMinimal}[thm]{1,2,8,6,15}}

\closeproofpartwide

\Needspace{8\baselineskip}

\proofpartwide{Fixpunktgleichheit}
  \setcounter{proofstepnr}{16}% Fortlaufende Schritte dieses Beweises.
\proofstepwidestar[1,2]{%
    \CBPart{F}{G}=(A\setminus G[B])\cup G[F[\CBPart{F}{G}]]}{%
    \FormulaRefAuto{A\subseteq B\dsep B\subseteq A\vdash A=B}[thm]{16,5}}

\closeproofpartwide
\end{tabproofsplitwide}
}
\Needspace{14\baselineskip}

\CSBComplementIdentity
\CSBProof{%
\begin{tabproofsplitwide}
  \Needspace{8\baselineskip}
  \hypertarget{csb.helper.CantorBernsteinComplementForward}{}
  \proofpartwideindR[Teilmengenrichtung \(\subseteq\)]{%
    F\colon A\inj B\dsep G\colon B\inj A\vdash
    G[B]\setminus G[F[\CBPart{F}{G}]]
      \subseteq A\setminus\CBPart{F}{G}
  }{%
    F\colon A\inj B\dsep G\colon B\inj A\vdash
    G[B]\setminus G[F[\CBPart{F}{G}]]
      \subseteq A\setminus\CBPart{F}{G}
  }[CantorBernsteinComplementForward]
    \proofstepwidestar[1]{F\colon A\inj B}{\rA}
    \proofstepwidestar[2]{G\colon B\inj A}{\rA}
    \proofstepwidestar[3]{y\in G[B]\setminus G[F[\CBPart{F}{G}]]}{\rA}
    \proofstepwidestar[3]{y\in G[B]}{%
      \FormulaRefAuto{x\in A\setminus B\vdash x\in A}[thm]{3}}
    \proofstepwidestar[3]{y\notin G[F[\CBPart{F}{G}]]}{%
      \FormulaRefAuto{x\in A\setminus B\vdash x\notin B}[thm]{3}}
    \proofstepwidestar[2]{G\colon B\to A}{%
      \FormulaRefAuto{Injektive Funktion}[def]{2}}
    \proofstepwidestar[2]{G[B]\subseteq A}{%
      \FormulaRefAuto{F\colon A\to B\vdash F[A]\subseteq B}[thm]{6}}
    \proofstepwidestar[2,3]{y\in A}{%
      \FormulaRefAuto{A\subseteq B\dsep x\in A\vdash x\in B}[thm]{7,4}}
  \proofstepwidestar[2,3]{( y \notin A \setminus { G [ B ] } \leftrightarrow y \in { G [ B ] } )}{%
      \FormulaRefAuto{RelCompNotMembership}[thm]{8}}
  \proofstepwidestar[2,3]{y\notin A\setminus G[B]}{%
      \FormulaRefAuto{P \leftrightarrow Q, Q \vdash P}[thm]{%
        9,4}}
    \proofstepwidestar[1,2]{%
      \CBPart{F}{G}=(A\setminus G[B])\cup G[F[\CBPart{F}{G}]]}{%
      \FormulaRefAuto{CantorBernsteinPartFixedPoint}[thm]{1,2}}
  \proofstepwidestar[1,2,3]{%
    y\notin(A\setminus G[B])\cup G[F[\CBPart{F}{G}]]}{%
      \FormulaRefAuto{NotInUnion}[thm]{10,5}}
    \proofstepwidestar[1,2,3]{y\notin\CBPart{F}{G}}{\rIE{11,12}}
    \proofstepwidestar[1,2,3]{y\in A\setminus\CBPart{F}{G}}{%
      \FormulaRefAuto{x\in A\setminus B\eqvdash x\in A\land x\notin B}[thm]{%
        \rAI{8,13}}}
    \proofstepwidestar[1,2]{%
      G[B]\setminus G[F[\CBPart{F}{G}]]
        \subseteq A\setminus\CBPart{F}{G}}{%
      \FormulaRefAuto{SubsetIntroductionRule}[thm]{[3]\vdots 14}}
  \closeproofpartwideindR

  \Needspace{8\baselineskip}

  \hypertarget{csb.helper.CantorBernsteinComplementBackward}{}
  \proofpartwideindR[Teilmengenrichtung \(\supseteq\)]{%
    F\colon A\inj B\dsep G\colon B\inj A\vdash
    A\setminus\CBPart{F}{G}
      \subseteq G[B]\setminus G[F[\CBPart{F}{G}]]
  }{%
    F\colon A\inj B\dsep G\colon B\inj A\vdash
    A\setminus\CBPart{F}{G}
      \subseteq G[B]\setminus G[F[\CBPart{F}{G}]]
  }[CantorBernsteinComplementBackward]
    \proofstepwidestar[1]{F\colon A\inj B}{\rA}
    \proofstepwidestar[2]{G\colon B\inj A}{\rA}
    \proofstepwidestar[3]{y\in A\setminus\CBPart{F}{G}}{\rA}
    \proofstepwidestar[3]{y\in A}{%
      \FormulaRefAuto{x\in A\setminus B\vdash x\in A}[thm]{3}}
    \proofstepwidestar[3]{y\notin\CBPart{F}{G}}{%
      \FormulaRefAuto{x\in A\setminus B\vdash x\notin B}[thm]{3}}
    \proofstepwidestar[1,2]{%
      \CBPart{F}{G}=(A\setminus G[B])\cup G[F[\CBPart{F}{G}]]}{%
      \FormulaRefAuto{CantorBernsteinPartFixedPoint}[thm]{1,2}}
    \proofstepwidestar[1,2,3]{%
      y\notin(A\setminus G[B])\cup G[F[\CBPart{F}{G}]]}{\rIE{6,5}}
  \proofstepwidestar[1,2,3]{y \notin { ( A \setminus G [ B ] ) } \land y \notin { G [ F [ \CBPart F G ] ] }}{%
      \FormulaRefAuto{OutsideUnionComponents}[thm]{7}}
  \proofstepwidestar[1,2,3]{y\notin A\setminus G[B]}{%
      \rAEa{8}}
  \proofstepwidestar[1,2,3]{y \notin { ( A \setminus G [ B ] ) } \land y \notin { G [ F [ \CBPart F G ] ] }}{%
      \FormulaRefAuto{OutsideUnionComponents}[thm]{7}}
  \proofstepwidestar[1,2,3]{y\notin G[F[\CBPart{F}{G}]]}{%
      \rAEb{10}}
    \proofstepwidestar[12]{y\notin G[B]}{\rA}
    \proofstepwidestar[3,12]{y\in A\setminus G[B]}{%
      \FormulaRefAuto{x\in A\setminus B\eqvdash x\in A\land x\notin B}[thm]{%
        \rAI{4,12}}}
    \proofstepwidestar[1,2,3,12]{\bot}{\rBI{13,9}}
    \proofstepwidestar[1,2,3]{\neg(y\notin G[B])}{\rCI{12,14}}
    \proofstepwidestar[1,2,3]{y\in G[B]}{\FormulaRefAuto{rDN}[thm]{15}}
    \proofstepwidestar[1,2,3]{%
      y\in G[B]\setminus G[F[\CBPart{F}{G}]]}{%
      \FormulaRefAuto{x\in A\setminus B\eqvdash x\in A\land x\notin B}[thm]{%
        \rAI{16,11}}}
    \proofstepwidestar[1,2]{%
      A\setminus\CBPart{F}{G}
        \subseteq G[B]\setminus G[F[\CBPart{F}{G}]]}{%
      \FormulaRefAuto{SubsetIntroductionRule}[thm]{[3]\vdots 17}}
  \closeproofpartwideindR

  \Needspace{8\baselineskip}

  \proofpartwide{Schluss}
    \proofstepwidestar[1]{F\colon A\inj B}{\rA}
    \proofstepwidestar[2]{G\colon B\inj A}{\rA}
    \proofstepwidestar[1,2]{%
      G[B]\setminus G[F[\CBPart{F}{G}]]
        \subseteq A\setminus\CBPart{F}{G}}{%
      \FormulaRefAuto{CantorBernsteinComplementForward}[thm]{1,2}}
    \proofstepwidestar[1,2]{%
      A\setminus\CBPart{F}{G}
        \subseteq G[B]\setminus G[F[\CBPart{F}{G}]]}{%
      \FormulaRefAuto{CantorBernsteinComplementBackward}[thm]{1,2}}
    \proofstepwidestar[1,2]{%
      G[B]\setminus G[F[\CBPart{F}{G}]]
        =A\setminus\CBPart{F}{G}}{%
      \FormulaRefAuto{A\subseteq B\dsep B\subseteq A\vdash A=B}[thm]{3,4}}
  \closeproofpartwide
\end{tabproofsplitwide}
}
\Needspace{14\baselineskip}

\CSBSecondBranchImage
\CSBProof{%
\begin{tabproofwide}
  \proofstepwidestar[1]{F\colon A\inj B}{\rA}
  \proofstepwidestar[2]{G\colon B\inj A}{\rA}
  \proofstepwidestar[1,2]{\CBPart{F}{G}\subseteq A}{%
    \FormulaRefAuto{CantorBernsteinPartSubset}[thm]{1,2}}
  \proofstepwidestar[1]{F\colon A\to B}{%
    \FormulaRefAuto{Injektive Funktion}[def]{1}}
  \proofstepwidestar[1,2]{F[\CBPart{F}{G}]\subseteq B}{%
    \FormulaRefAuto{C\subseteq A\dsep F\colon A\to B
      \vdash F[C]\subseteq B}[thm]{3,4}}
  \proofstepwidestar[1,2]{B \subseteq B}{%
      \FormulaRefAuto{A\subseteq A}[thm]}
  \proofstepwidestar[1,2]{%
    G[B\setminus F[\CBPart{F}{G}]]
      =G[B]\setminus G[F[\CBPart{F}{G}]]}{%
    \FormulaRefAuto{InjectiveImageDifference}[thm]{2,5,
      6}}
  \proofstepwidestar[1,2]{%
    G[B]\setminus G[F[\CBPart{F}{G}]]
      =A\setminus\CBPart{F}{G}}{%
    \FormulaRefAuto{CantorBernsteinComplementIdentity}[thm]{1,2}}
  \proofstepwidestar[1,2]{%
    G[B\setminus F[\CBPart{F}{G}]]
      =A\setminus\CBPart{F}{G}}{\rChain{7,8}}
\end{tabproofwide}
}
\Needspace{14\baselineskip}

\CSBFixedInjections
\CSBProof{%
\begin{tabproofsplitwide}
\Needspace{8\baselineskip}
\proofpartwide{Erster bijektiver Zweig}
\proofstepwidestar[1]{F\colon A\inj B}{\rA}
  \proofstepwidestar[2]{G\colon B\inj A}{\rA}
  \proofstepwidestar[1,2]{\CBPart{F}{G}\subseteq A}{%
    \FormulaRefAuto{CantorBernsteinPartSubset}[thm]{1,2}}
  \proofstepwidestar[1,2]{%
    F\restriction_{\CBPart{F}{G}}\colon
      \CBPart{F}{G}\bij F[\CBPart{F}{G}]}{%
    \FormulaRefAuto{InjectionRestrictionToImageBijective}[thm]{1,3}}

\closeproofpartwide

\Needspace{8\baselineskip}

\proofpartwide{Zweiter bijektiver Zweig}
  \setcounter{proofstepnr}{4}% Fortlaufende Schritte dieses Beweises.
\proofstepwidestar[1]{F\colon A\to B}{%
    \FormulaRefAuto{Injektive Funktion}[def]{1}}
  \proofstepwidestar[1,2]{F[\CBPart{F}{G}]\subseteq B}{%
    \FormulaRefAuto{C\subseteq A\dsep F\colon A\to B
      \vdash F[C]\subseteq B}[thm]{3,5}}
  \proofstepwidestar[1,2]{B\setminus F[\CBPart{F}{G}]\subseteq B}{%
    \FormulaRefAuto{A\setminus B\subseteq A}[thm]}
  \proofstepwidestar[1,2]{%
    G\restriction_{B\setminus F[\CBPart{F}{G}]}\colon
      B\setminus F[\CBPart{F}{G}]
      \bij G[B\setminus F[\CBPart{F}{G}]]}{%
    \FormulaRefAuto{InjectionRestrictionToImageBijective}[thm]{2,7}}
  \proofstepwidestar[1,2]{%
    G[B\setminus F[\CBPart{F}{G}]]
      =A\setminus\CBPart{F}{G}}{%
    \FormulaRefAuto{CantorBernsteinSecondBranchImage}[thm]{1,2}}
  \proofstepwidestar[1,2]{%
    G\restriction_{B\setminus F[\CBPart{F}{G}]}\colon
      B\setminus F[\CBPart{F}{G}]
      \bij A\setminus\CBPart{F}{G}}{\rIE{9,8}}
  \proofstepwidestar[1,2]{%
    \bigl(G\restriction_{B\setminus F[\CBPart{F}{G}]}\bigr)^{-1}
      \colon A\setminus\CBPart{F}{G}
      \bij B\setminus F[\CBPart{F}{G}]}{%
    \FormulaRefAuto{InverseFunctionBijection}[thm]{10}}

\closeproofpartwide

\Needspace{8\baselineskip}

\proofpartwide{Disjunktheit und Verkleben}
  \setcounter{proofstepnr}{11}% Fortlaufende Schritte dieses Beweises.
\proofstepwidestar[]{A\setminus\CBPart{F}{G}\subseteq A\setminus\CBPart{F}{G}}{%
      \FormulaRefAuto{A\subseteq A}[thm]}
\proofstepwidestar[]{%
    \CBPart{F}{G}\cap(A\setminus\CBPart{F}{G})=\varnothing}{%
    \FormulaRefAuto{K\subseteq H\setminus M\vdash M\cap K=\varnothing}[thm]{%
      12}}
  \proofstepwidestar[]{B\setminus F[\CBPart{F}{G}]\subseteq B\setminus F[\CBPart{F}{G}]}{%
      \FormulaRefAuto{A\subseteq A}[thm]}
  \proofstepwidestar[]{%
    F[\CBPart{F}{G}]\cap
      (B\setminus F[\CBPart{F}{G}])=\varnothing}{%
    \FormulaRefAuto{K\subseteq H\setminus M\vdash M\cap K=\varnothing}[thm]{%
      14}}
  \proofstepwidestar[1,2]{%
    \CaseFun{\CBPart{F}{G}}{A\setminus\CBPart{F}{G}}
      {F\restriction_{\CBPart{F}{G}}}
      {\bigl(G\restriction_{B\setminus F[\CBPart{F}{G}]}\bigr)^{-1}}
      \colon
      \CBPart{F}{G}\cup(A\setminus\CBPart{F}{G})
      \bij F[\CBPart{F}{G}]
        \cup(B\setminus F[\CBPart{F}{G}])}{%
    \FormulaRefAuto{CaseFunBijectiveDisjointTargets}[thm]{13,15,4,11}}

\closeproofpartwide

\Needspace{8\baselineskip}

\proofpartwide{Identifikation von Definitions- und Zielbereich}
  \setcounter{proofstepnr}{16}% Fortlaufende Schritte dieses Beweises.
\proofstepwidestar[1,2]{%
    A=\CBPart{F}{G}\cup(A\setminus\CBPart{F}{G})}{%
    \FormulaRefAuto{M\subseteq H\vdash H=M\cup(H\setminus M)}[thm]{3}}
  \proofstepwidestar[1,2]{%
    B=F[\CBPart{F}{G}]\cup(B\setminus F[\CBPart{F}{G}])}{%
    \FormulaRefAuto{M\subseteq H\vdash H=M\cup(H\setminus M)}[thm]{6}}
  \proofstepwidestar[1,2]{%
    \CaseFun{\CBPart{F}{G}}{A\setminus\CBPart{F}{G}}
      {F\restriction_{\CBPart{F}{G}}}
      {\bigl(G\restriction_{B\setminus F[\CBPart{F}{G}]}\bigr)^{-1}}
      \colon A\bij B}{\rIE{17,\rIE{18,16}}}
  \proofstepwidestar[1,2]{\exists H\colon A\bij B}{\rEI{19}}

\closeproofpartwide
\end{tabproofsplitwide}
}

\fi
\fi

\end{document}
